\subsubsection{Prikačivanje na pokrenut proces}\label{gdb-attach-process}

U realnom radu problem se često pojavi u procesu koji je već pokrenut. \texttt{gdb} može da se prikači na takav proces, zaustavi ga, pročita stanje i zatim se odvoji bez namernog gašenja procesa.

Primer simulira mali servis koji periodično obrađuje zahteve:

\codefile{service.c}{c}

U prvom terminalu pokrenemo servis:

\begin{minted}{bash}
$ make
$ ./service
pid=12345
\end{minted}

U drugom terminalu se prikačimo na PID koji je program ispisao:

\begin{minted}{bash}
$ gdb -p 12345
\end{minted}

Proces se zaustavlja dok je \texttt{gdb} prikačen. Tada možemo pročitati stek i globalno stanje:

\begin{minted}{text}
(gdb) bt
(gdb) p requests
(gdb) p slow_path
(gdb) info threads
(gdb) continue
^C
(gdb) p requests
(gdb) detach
(gdb) quit
\end{minted}

Ovaj tok je koristan za procese koji troše CPU, vise u čekanju ili imaju stanje koje je teško reprodukovati od početka. Ipak, prikačivanje nije neutralno: proces je zaustavljen dok ga pregledamo, što može izazvati timeout-e ili promeniti ponašanje sistema.

Na Linux sistemima prikačivanje može biti ograničeno podešavanjem \path{/proc/sys/kernel/yama/ptrace_scope}. Ako se dobije greška \texttt{Operation not permitted}, proveriti vlasništvo procesa, capability-je i to podešavanje. Na produkcionim sistemima odluka da se prikači debager mora biti deo operativne procedure, jer debager ima pun pristup memoriji procesa.
